\documentclass[12pt,a4paper]{article}

\usepackage[utf8]{inputenc}
\usepackage[T1]{fontenc}
\usepackage[ngerman]{babel}
\usepackage{geometry}
\usepackage{hyperref}
\usepackage{setspace}

\geometry{a4paper, margin=2.5cm}
\onehalfspacing

\title{Konzeption und Implementierung eines Python-basierten Planungstools zur Unterstützung der Lehrveranstaltungsplanung}
\author{Nicolas}
\date{\today}

\begin{document}

\maketitle

\tableofcontents
\newpage

\section{Einleitung}

\subsection{Problemstellung}

Es werden jährlich neue Lehrveranstaltungen, Übungen, Labore und Klausuren eingeplant. Dazu müssen bei den Entscheidungen, wann welcher Termin eingeplant wird, mehrere verscheidene Informationen gleichzeitig berücksichtigt werden:

\begin{itemize}
\item zeitliche Überschneidungen von Lehrveranstaltungen, Lehrpersonen, Räumen und Gruppen
\item Raumkapazitäten
\item freie Tage/Feiertage
\end{itemize}

Diese Planung wird durch die hohe Menge an individuellen Termine eine komplexe Aufgabe. Die manuelle Planung erfolgt bisher über eine Zusammensetzung an Excel-Tabellen, welche die Verschiebung oder Veränderung einzelner Termine nicht effizient erlauben.
Durch die manuelle Planung können auch weitere Probleme auftreten wie doppelte oder uneinheitliche Daten, Änderungen müssen an mehreren Stellen angepasst werden, schwer erkennbare Konflikte und ein generell schlechter Überblick über die einzelnen Termine. Beispiele für mögliche Konflikte könnten sein: doppelt belegte Räume, Lehrperson zeitgleich eingeplant oder ein Termin an einem Feiertag angesetzt.

Auf Grund dieser Probleme und dadurch dass dies eine jährlich wiederkehrende Aufgabe ist besteht ein Bedarf an einer zentralen Softwarelösung. Diese Software soll die Planung der Termine übersichtlicher, strukturierter und nachvollziehbarer machen.
Daher geht es in dieser Arbeit um die Entwicklung eines Python-basierten Planungstool für die Unterstützung der LVA-Jahresplanung. Der Fokus liegt dabei auf der praktischen Anwendbarkeit dieser Software am Institut.

\subsection{Zielsetzung}

Ziel dieser Bachlorarbeit ist die Konzeption, Implementierung und Bewertung eines Planungstools. Dieses Tool soll den jährlichen Planungprozess der LVA-Termine digital unterstützen und damit effizienter und einfacher gestalten.
Zentrale Ziele dieser Software sind:

\begin{itemize}
\item strukturierte und zentrale Verwaltung von Planungsdaten
\item übersichtliche Darstellung von Terminen
\item einfaches Verschieben mittels Drag-and-Drop und Zuweisen von Terminen
\item bessere Erkennung von Planungskonflikten
\item Filterung nach relevanten Kriterien
\item Import und Export von Planungsdaten
\end{itemize}

Die Umsetzung soll in Form einer mit Python entwickelten Desktop-Anwendung erfolgen. Die Entwicklung einer einfachen, intuitiven grafischen Benutzeroberfläche ist dabei besonders wichtig, da die Bedienung dieses Tools auch ohme Programmierkentnisse möglich sein soll.

Ziel dieser Anwendung ist nicht die gesamte Planung zu übernehmen, sondern es soll als Unterstützung für den User dienen, der die Planung durführt. Es soll besonders auf benutzerfreundlichkeit, nachvollziehbarkeit, erweiterbarkeit und praktische Nutzbarkeit geachtet werden.


\subsection{Forschungsfrage}

Aus der beschriebenen Problemstellung und Zielsetzung der Arbeit und dem Bedarf nach einer digitalen Unterstützung für die LVA-Jahresplanung ergibt sich die Frage, wie man ein Werkzeug entwickelt, das Planungsdaten strukturiert verwaltet, Termine übersichtlich darstellt und mögliche Konflikte leichter erkennbar macht.

Die zentrale Forschungsfrage dieser Arbeit lautet daher:

\begin{quote}
Wie kann ein intuitives Python-basiertes Desktop-Planungstool entwickelt werden, um die jährliche Lehrveranstaltungsplanung an einem Institut übersichtlicher, strukturierter, effizienter und weniger fehleranfällig zu gestalten?
\end{quote}

Um diese Frage zu beantworten werden auch zusätzliche Teilfragen betrachtet:

\begin{itemize}
    \item Welche Anforderungen ergeben sich aus dem Planungsprozess?
    \item Wie können Termine, Lehrveranstaltungen, Räume und weitere Daten sinnvoll strukturiert und verwaltet werden?
    \item Wie kann eine Benutzeroberfläche gestalten, sodass Termine übersichtlich dargestellt, gefiltert und mittels Drag-and-Drop verschoben werden können?
    \item Wie können typische Planungskonflikte erkannt und für den User sichtbar schnell gemacht werden?
    \item Inwiefern erfüllt das entwickelte Tool die definierten Anforderungen?
\end{itemize}

\subsection{Aufbau der Arbeit}

Die Arbeit beginnt mit der Beschreibung der Probleme, die durch die jährliche Planung der Lehrveranstaltungen entstehen. Danach werden die Ziele der Arbeit und die Forschungsfrage vorgestellt. Dadurch wird deutlich gemacht, warum diese Anwendung entwickelt wurde und welche Punkte in der Arbeit besonders wichtig sind.

...

\newpage

\section{Grundlagen und Kontext}

\subsection{Lehrveranstaltungsplanung an einem Institut}

Bei der Lehrveranstaltungsplanung geht es darum, die verschiedenen Lehrveranstaltungen eines Instituts in einen passenden zeitlichen und organisatorischen Rahmen zu bringen. Dazu müssen Veranstaltungen wie Vorlesungen, Übungen, Labore oder Prüfungen konkreten Terminen, Räumen und teilweise auch bestimmten Gruppen zugeordnet werden. Die Planung ist daher mehr als nur das Eintragen von Terminen in einen Kalender.

Je nach Art der Lehrveranstaltung können dabei unterschiedliche Anforderungen entstehen. Eine Vorlesung benötigt zum Beispiel oft einen größeren Raum, während ein Labor an bestimmte Ausstattung gebunden sein kann. Prüfungen wiederum finden häufig nur an einzelnen Terminen statt und müssen besonders klar geplant werden. Auch die Dauer der Termine und die Anzahl der teilnehmenden Personen können sich je nach Veranstaltung unterscheiden.

Ein Termin steht dabei meistens nicht für sich allein, sondern ist mit mehreren anderen Informationen verbunden. Dazu gehören zum Beispiel die Lehrveranstaltung selbst, das Semester, die zuständige Lehrperson, der Raum oder die betroffene Studierendengruppe. Dadurch entsteht eine Struktur, bei der viele Daten miteinander zusammenhängen.

Für ein Institut ist es wichtig, dass aus diesen Informationen ein übersichtlicher und nachvollziehbarer Plan entsteht. Dieser Plan soll sowohl für die Verwaltung als auch für Lehrende und Studierende verständlich sein.

In dieser Arbeit wird die Lehrveranstaltungsplanung deshalb als ein Prozess verstanden, bei dem verschiedene Planungsdaten gesammelt, miteinander verknüpft und für die weitere Bearbeitung nutzbar gemacht werden. Diese Sichtweise bildet die Grundlage für die spätere Entwicklung des Planungstools.

\subsection{Herausforderungen manueller Planungsprozesse}

Bei der bisherigen Planung werden viele Informationen in Tabellen gesammelt. Das ist am Anfang praktisch, weil Tabellen schnell erstellt und angepasst werden können. Wenn jedoch viele Termine eingeplant werden müssen, wird es mit der Zeit schwer, den Überblick zu behalten. Besonders bei der Jahresplanung kommen viele einzelne Termine zusammen, die sich gegenseitig beeinflussen können.

Ein Problem ist, dass Änderungen nicht immer nur eine Stelle betreffen. Wenn zum Beispiel ein Termin verschoben wird, muss geprüft werden, ob der neue Zeitpunkt mit anderen Terminen zusammenpasst. Außerdem muss kontrolliert werden, ob der Raum frei ist und ob die betroffene Lehrperson zu diesem Zeitpunkt nicht bereits anders eingeplant ist. Diese Kontrolle muss bei einer manuellen Planung selbst gemacht werden.

Dadurch wird die manuelle Planung mit zunehmender Anzahl an Terminen immer aufwendiger. Tabellen bleiben zwar flexibel, bieten aber nur begrenzte Unterstützung dabei, Termine übersichtlich darzustellen und mögliche Konflikte frühzeitig zu erkennen.

\subsection{Softwaregestützte Planung}

Eine softwaregestützte Planung kann den Planungsprozess vor allem dabei unterstützen, viele einzelne Informationen besser zusammenzuführen. Anstatt Termine nur in Tabellen einzutragen, können sie in einer Anwendung direkt mit weiteren Daten verbunden werden, zum Beispiel mit einer Lehrveranstaltung, einem Raum oder einer Lehrperson. Dadurch wird die Planung nicht nur gesammelt, sondern auch besser nutzbar gemacht.

Ein Vorteil ist, dass die Daten an einer zentralen Stelle verwaltet werden können. Wenn ein Termin geändert wird, muss diese Änderung nicht in mehreren Tabellen oder Dateien nachgetragen werden. Die angepassten Daten können direkt in der Anwendung angezeigt und weiterverwendet werden.

Auch die Darstellung spielt eine wichtige Rolle. Eine Kalenderansicht ist für die Planung oft verständlicher als eine reine Tabelle, weil Termine zeitlich eingeordnet werden können. So sieht man schneller, an welchen Tagen bereits viele Termine liegen oder wo noch freie Zeiträume vorhanden sind. Durch Filter kann zusätzlich gezielt nach bestimmten Lehrveranstaltungen, Räumen, Lehrpersonen, etc. gesucht werden.

Besonders hilfreich ist außerdem eine direkte Bearbeitung der Termine über die Benutzeroberfläche. Wenn Termine zum Beispiel per Drag-and-Drop verschoben oder zugewiesen werden, wird die Planung übersichtlich und schnell angepasst. Statt Werte manuell in Tabellenzellen zu ändern, kann der Nutzer einen Termin direkt an die passende Stelle ziehen. Dadurch wird der Planungsprozess näher an der tatsächlichen Kalenderansicht durchgeführt.

Eine Software kann außerdem dabei helfen, mögliche Konflikte sichtbarer zu machen. Wenn ein Raum bereits belegt ist oder eine Lehrperson zur gleichen Zeit in einem anderen Termin vorkommt, kann das Programm solche Fälle erkennen und anzeigen. Die Entscheidung, wie mit einem Konflikt umgegangen wird, bleibt weiterhin beim Nutzer. Die Software übernimmt also nicht die gesamte Planung, sondern unterstützt dabei, Fehler schneller zu finden und Änderungen übersichtlicher umzusetzen.


\subsection{Eingesetzte Technologien}

Für die Umsetzung des Planungstools wurden mehrere Technologien verwendet. Die Anwendung wurde mit Python entwickelt und als Desktop-Programm umgesetzt. Für die grafische Oberfläche kommt PySide6 zum Einsatz, da damit die verschiedenen Fenster, Tabellen und Bedienbereiche der Anwendung erstellt werden können.

Die Daten des Planungstools werden in JSON-Dateien gespeichert. Dadurch können die einzelnen Planungsdaten strukturiert abgelegt und beim Start der Anwendung wieder geladen werden. Außerdem eignet sich dieses Format gut für den Import und Export von Projektdaten.

Die wichtigsten verwendeten Technologien werden in den folgenden Abschnitten kurz beschrieben.

\subsubsection{Python}

Das Planungstool wurde mit Python entwickelt.\footnote{\url{https://www.python.org/}} Python bildet dabei die Grundlage für die eigentliche Logik der Anwendung, also zum Beispiel für das Laden und Speichern der Daten, das Bearbeiten von Terminen und das Prüfen von möglichen Konflikten.

Für dieses Projekt war Python passend, weil die Sprache gut lesbar ist und sich der Code übersichtlich strukturieren lässt. Ein weiterer Grund für die Verwendung von Python ist, dass die Sprache weit verbreitet ist und von vielen Personen verwendet wird. Dadurch kann die spätere Wartung oder Erweiterung des Tools erleichtert werden, da sich neue Entwicklerinnen und Entwickler schneller in den Code einarbeiten können.

Außerdem unterstützt Python viele Bibliotheken. Für dieses Projekt war vor allem PySide6 wichtig, weil damit die grafische Oberfläche erstellt wurde. Dadurch konnten die verschiedenen Fenster, Tabellen und Bedienelemente des Planungstools umgesetzt werden.

\subsubsection{PySide6 und Qt}

Für die grafische Benutzeroberfläche des Planungstools wurde PySide6 verwendet. PySide6 verbindet Python mit dem Qt-Framework. Qt bietet viele Elemente, die für Desktop-Anwendungen gebraucht werden, zum Beispiel Fenster, Tabellen, Buttons, Menüs und Dialoge.\footnote{\url{https://doc.qt.io/qt-6/}} Dadurch konnte das Tool als lokale Anwendung mit einer eigenen Benutzeroberfläche umgesetzt werden.

Grundsätzlich hätte die Oberfläche auch mit PyQt entwickelt werden können, da PyQt und PySide6 ähnliche Funktionen bereitstellen und beide auf Qt basieren. Für dieses Projekt wurde jedoch PySide6 gewählt, da es das offizielle Qt-for-Python-Binding ist.\footnote{\url{https://doc.qt.io/qtforpython-6/index.html}} Außerdem kann die LGPL-Lizenz von PySide6 bei eigenständigen Anwendungen flexibler sein.\footnote{\url{https://www.qt.io/development/open-source-lgpl-obligations}}

Im Planungstool wird PySide6 vor allem für die verschiedenen Bedienbereiche der Anwendung genutzt. Dazu gehören zum Beispiel Kalenderansichten, Tabellen, Eingabefenster, Filterbereiche und Listen.

Wenn der Nutzer einen Termin auswählt, einen Filter setzt oder Daten in einem Dialog bearbeitet, muss die Anwendung auf diese Eingaben reagieren. PySide6 stellt dafür die Verbindung zwischen der Benutzerberfläche und der Programmlogik her. Dadurch können Änderungen direkt verarbeitet und an anderen Stellen der Anwendung sichtbar gemacht werden.

\subsubsection{JSON als Datenformat}

Die Planungsdaten werden in JSON-Dateien gespeichert.\footnote{\url{https://www.rfc-editor.org/rfc/rfc8259}} Dazu gehören zum Beispiel Termine, Lehrveranstaltungen, Räume, Semester und weitere Informationen, die für die Planung benötigt werden. JSON wurde verwendet, weil sich damit strukturierte Daten einfach speichern, laden und bei Bedarf manuell bearbeiten lassen.

Da die JSON-Dateien auch außerhalb des Programms geöffnet werden können, bleibt nachvollziehbar, welche Daten gespeichert wurden. Das kann zum Beispiel hilfreich sein, wenn Fehler gesucht werden oder überprüft werden soll, ob Daten richtig übernommen wurden.

Grundsätzlich hätten für die Speicherung auch CSV-Dateien verwendet werden können. CSV eignet sich vor allem für einfache tabellarische Daten. Für das Planungstool ist das jedoch nicht immer passend, da die Daten teilweise stärker miteinander verknüpft sind. Ein Termin kann zum Beispiel mehrere Eigenschaften oder Verweise auf andere Daten enthalten, wie auf eine LVA oder einen Raum.
JSON ist dafür besser geeignet, weil zusammenhängende Informationen als Objekte gespeichert werden können. Dadurch lassen sich die Planungsdaten strukturierter abbilden als in einer einfachen Tabelle.

Für die aktuelle Version und Größe der Anwendung ist diese Art der Speicherung ausreichend. Es ist also kein zentraler Server oder eine Datenbank notwendig. Wenn das Tool später erweitert wird, könnte eine Datenbank jedoch eine sinnvolle Weiterentwicklung sein.

\newpage

\section{Anforderungsanalyse}

\subsection{Ableitung der Anforderungen}

Die Anforderungen an das Planungstool ergeben sich aus den typischen Arbeitsschritten der jährlichen Lehrveranstaltungsplanung. Dabei müssen Termine nicht nur gespeichert, sondern auch mit Informationen wie Lehrveranstaltungen, Räumen, Semestern oder Gruppen verbunden werden.

Für die Anforderungsanalyse wurden daher vor allem das Anlegen und Bearbeiten von Terminen, das Verschieben von Terminen, das Filtern der Daten und das Erkennen möglicher Konflikte betrachtet. Daraus ergeben sich die folgenden funktionalen und nicht-funktionalen Anforderungen.

\subsection{Funktionale Anforderungen}

Funktionale Anforderungen beschreiben, welche Aufgaben die Software erfüllen soll. Die wichtigsten funktionalen Anforderungen des Planungstools sind in Tabelle~\ref{tab:functional-requirements} dargestellt.

\begin{table}[h]
\centering
\begin{tabular}{|p{4cm}|p{9cm}|}
\hline
\textbf{Anforderung} & \textbf{Beschreibung} \\
\hline
Terminverwaltung & Termine sollen erstellt, bearbeitet und gelöscht werden können. \\
\hline
Datenverwaltung & Lehrveranstaltungen, Räume, Semester, Fachrichtungen und weitere Daten sollen verwaltet werden können. \\
\hline
Terminzuordnung & Termine sollen mit Lehrveranstaltungen, Räumen, Semestern oder Gruppen verbunden werden können. \\
\hline
Kalenderdarstellung & Termine sollen in einer übersichtlichen Kalenderansicht dargestellt werden. \\
\hline
Drag-and-Drop & Termine sollen direkt über die Benutzeroberfläche verschoben oder zugewiesen werden können. \\
\hline
Filterung & Termine sollen nach bestimmten Kriterien gefiltert werden können. \\
\hline
Konflikterkennung & Mögliche Konflikte, wie doppelt belegte Räume oder zeitliche Überschneidungen, sollen erkannt und angezeigt werden. \\
\hline
Speicherung & Planungsdaten sollen gespeichert und beim erneuten Start wieder geladen werden können. \\
\hline
Import und Export & Planungsdaten sollen importiert und exportiert werden können. \\
\hline
\end{tabular}
\caption{Funktionale Anforderungen an das Planungstool}
\label{tab:functional-requirements}
\end{table}

\subsection{Nicht-funktionale Anforderungen}

Nicht-funktionale Anforderungen beschreiben, welche Eigenschaften die Software haben soll. Sie beziehen sich also nicht direkt auf einzelne Funktionen, sondern auf Aspekte wie Bedienbarkeit, Wartbarkeit oder Stabilität.

\begin{table}[h]
\centering
\begin{tabular}{|p{4cm}|p{9cm}|}
\hline
\textbf{Anforderung} & \textbf{Beschreibung} \\
\hline
Benutzerfreundlichkeit & Das Tool soll auch ohne Programmierkenntnisse einfach bedienbar sein. \\
\hline
Übersichtlichkeit & Termine und Planungsdaten sollen klar dargestellt werden. \\
\hline
Wartbarkeit & Der Code soll übersichtlich aufgebaut sein, damit Fehler behoben und bestehende Funktionen angepasst werden können. \\
\hline
Erweiterbarkeit & Neue Funktionen sollen später ergänzt werden können. \\
\hline
Nachvollziehbarkeit & Daten und Änderungen sollen verständlich und kontrollierbar bleiben. \\
\hline
Stabilität & Die Anwendung soll Daten zuverlässig speichern und laden können. \\
\hline
Lokale Nutzbarkeit & Das Tool soll direkt auf einem Computer ausgeführt werden können und keinen Server benötigen. \\
\hline
\end{tabular}
\caption{Nicht-funktionale Anforderungen an das Planungstool}
\label{tab:nonfunctional-requirements}
\end{table}

\subsection{Abgrenzung der Software}

Das Planungstool wurde als Unterstützung für die jährliche Lehrveranstaltungsplanung entwickelt. Es soll den Planungsprozess übersichtlicher machen und bei der Bearbeitung der Termine helfen. Die organisatorischen Entscheidungen werden dabei aber alle weiterhin vom Nutzer getroffen.

Nicht Teil dieser Version ist eine automatische Optimierung der gesamten Planung. Die Software erstellt also nicht selbstständig den bestmöglichen Stundenplan, sondern unterstützt den bestehenden Planungsprozess. Auch eine direkte Anbindung an offizielle Universitätssysteme oder eine automatische Raumreservierung ist aktuell nicht umgesetzt.

Die Anwendung ist außerdem als lokale Desktop-Anwendung ausgelegt. Ein gleichzeitiger Zugriff durch mehrere Nutzer ist in dieser Version nicht vorgesehen. Funktionen wie eine Anbindung an bestehende Systeme, automatische Raumreservierung oder Mehrnutzerbetrieb könnten sinnvolle Erweiterungen für eine spätere Version des Tools sein.

\newpage

\section{Konzeption des Planungstools}

\subsection{Gesamtarchitektur der Anwendung}

Das Planungstool wurde als lokale Desktop-Anwendung konzipiert und besteht nicht nur aus einer einzelnen Kalenderansicht, sondern aus mehreren Bereichen, die zusammenarbeiten. Im Mittelpunkt steht das Hauptfenster(Main Window im Code), das die verschiedenen Teile der Benutzeroberfläche verbindet und die zentralen Abläufe der Anwendung steuert.

Grundsätzlich lässt sich die Anwendung in drei Bereiche einteilen: die Benutzeroberfläche, die Datenverwaltung und die Planungslogik. Die Benutzeroberfläche stellt Kalenderansichten, Listen, Eingabefenster und weitere Bedienelemente bereit. Die Datenverwaltung ist dafür zuständig, Planungsdaten zu laden, zu speichern und für die anderen Teile der Anwendung verfügbar zu machen. Die Planungslogik umfasst Funktionen wie Filterung, Konflikterkennung und die Aktualisierung der Anzeige nach Änderungen.

\subsection{Aufbau der Benutzeroberfläche}

Die Benutzeroberfläche ist in mehrere Bereiche aufgeteilt. Der zentrale Bereich ist der Planner Workspace, in dem die Termine in unterschiedlichen Kalenderansichten dargestellt werden. Es eine Tagesansicht, eine Wochenansicht und eine Monatsansicht. Die Tagesansicht eignet sich besonders für eine genauere Planung nach Uhrzeit und Raum. Die Wochenansicht gibt einen Überblick über mehrere Tage, während die Monatsansicht eher zur groben zeitlichen Orientierung dient.

Neben den Kalenderansichten gibt es weitere Bereiche, die als sogenannte Docks umgesetzt sind.\footnote{\url{https://doc.qt.io/qt-6/qdockwidget.html}} Dazu gehören die Terminliste, der Dateneditor, die Datumsnavigation, der globale Filterbereich und die Konfliktanzeige. Durch die verschiebbaren Docks kann jeder Benutzer seinen Arbeitsbereich so einrichten, wie es für ihn am besten passt. Das gewählte Layout lässt sich speichern und später wiederverwenden.

Die Terminliste dient dazu, vorhandene Termine übersichtlich anzuzeigen. Sie ist besonders wichtig für Termine, die noch nicht vollständig eingeplant wurden. Der Dateneditor wird verwendet, um Termine und Stammdaten wie Lehrveranstaltungen, Räume, Semester, Fachrichtungen oder freie Tage zu erstellen, zu bearbeiten oder zu löschen.

Der Filterbereich grenzt die Terminliste gezielt ein, zum Beispiel nach Lehrveranstaltung, Raum, Dozent oder Semester. Die Konfliktanzeige zeigt eine interaktive Liste aller aktuellen Planungskonflikte. Mit dem Navigationsbereich kann der Benutzer beliebig zwischen den Ansichten wechseln und innerhalb der jeweiligen Zeitintervalle navigieren.

\subsection{Datenmodell und Datenbeziehungen}

Das Datenmodell bildet die Grundlage des Planungstools. Im Mittelpunkt stehen die Termine, da sie die eigentlichen Planungseinträge darstellen. Ein Termin besteht aber nicht nur aus einem Datum und einer Uhrzeit, sondern ist mit weiteren Informationen verbunden. Dazu gehören zum Beispiel die zugehörige Lehrveranstaltung, ein Raum, ein Semester, ein Veranstaltungstyp oder eine Gruppe.

Neben den Terminen gibt es noch andere Datenbereiche. Dazu zählen Lehrveranstaltungen, Räume, Semester, geplante Semester, Fachrichtungen und freie Tage. Diese Daten werden getrennt gespeichert, damit sie nicht bei jedem Termin erneut eingetragen werden müssen. Stammdaten können dadurch einmal angelegt und anschließend mehrfach verwendet werden.

Die Verbindung zwischen den Daten erfolgt über IDs. Ein Termin speichert also nicht alle Informationen direkt ab, sondern verweist auf den passenden Eintrag. Dadurch können Änderungen an Stammdaten zentral vorgenommen werden.

Neben den Planungsdaten speichert die Anwendung auch Einstellungen und Konfliktpräferenzen in JSON-Dateien. So bleiben persönliche Anpassungen, etwa gespeicherte Layouts oder ausgewählte Konfliktregeln, auch beim nächsten Start erhalten.

\subsection{Planungsworkflow}

Der Arbeitsablauf im Planungstool beginnt in der Regel mit dem Laden der vorhandenen Daten. Danach können bei Bedarf Stammdaten ergänzt oder angepasst werden. Auf dieser Grundlage werden die einzelnen Termine erstellt und weiter geplant.

Ein Termin kann entweder direkt mit Datum, Uhrzeit und Raum angelegt werden oder zunächst ohne feste Zuordnung bestehen. Noch nicht zugewiesene Termine können später aus der Terminliste in die Kalenderansicht per Drag-and-Drop übernommen werden. Während des Drag-and-Drops werden mögliche Konflikte direkt als rote Vorschau im Kalender angezeigt.

Während der Planung können Filter genutzt werden, um nur bestimmte Termine anzuzeigen. Die Konfliktanzeige hilft zusätzlich dabei, mögliche Probleme schnell zu erkennen.

Der Workflow ist so aufgebaut, dass der Nutzer die Planung weiterhin selbst steuert. Das Tool zeigt Termine und Konflikte übersichtlich an und erleichtert Änderungen.

\subsection{Aktualisierung und Zusammenspiel der Komponenten}

Da das Planungstool aus mehreren Bereichen besteht, müssen diese nach Änderungen miteinander abgestimmt werden. Wenn ein Termin erstellt, bearbeitet oder verschoben wird, soll diese Änderung nicht nur an einer Stelle sichtbar sein. Auch die Kalenderansicht, die Terminliste und die Konfliktanzeige müssen den neuen Stand übernehmen.

Dafür ist das Hauptfenster als zentrale Stelle der Anwendung wichtig. Es verbindet die verschiedenen Bereiche miteinander und sorgt dafür, dass nach relevanten Änderungen eine Aktualisierung ausgelöst wird. Die einzelnen Komponenten laden anschließend ihre Inhalte neu oder passen ihre Darstellung an.

\newpage

\section{Implementierung}

\subsection{Entwicklungsumgebung und Projektstruktur}

Die Anwendung wurde mit Python in der Version 3.13.2 entwickelt.\footnote{\url{https://www.python.org/downloads/release/python-3132/}} Für die grafische Benutzeroberfläche wird PySide6 verwendet. Die Implementierung des Planungstools erfolgte in Visual Studio Code.

Der Quellcode ist in mehrere Bereiche aufgeteilt. Im \texttt{src}-Verzeichnis befinden sich die wichtigsten Bestandteile des Programms. Dazu gehören die Benutzeroberfläche, verschiedene Services, Dialoge, Hilfsfunktionen und die zentrale Programmlogik. Gestartet wird die Anwendung über die Datei \texttt{main.py}.

Die Planungsdaten liegen in JSON-Dateien und befinden sich je nach ausgewähltem Datenpfad zum Beispiel im Ordner \texttt{data}. Dadurch kann das Tool mit unterschiedlichen Datensätzen verwendet werden.

Für die benötigten Pakete wird eine \texttt{requirements.txt}-Datei verwendet. Zusätzlich enthält das Projekt eine \texttt{pyproject.toml}, in der weitere Projektinformationen und Einstellungen festgelegt werden.

\subsection{Hauptfenster und zentrale Steuerung}

Das Hauptfenster der Anwendung wird durch die Klasse \texttt{MainWindow} umgesetzt. In diesem Fenster kommen die wichtigsten Teile des Planungstools zusammen. Dazu gehören der Kalenderbereich, die Terminliste, der Dateneditor, die Filter, die Konfliktanzeige und die Datumsnavigation.

Das folgende Beispiel zeigt vereinfacht, wie diese Bereiche im Hauptfenster eingebunden werden:

\begin{verbatim}
self.planner = PlannerWorkspace(self, self.ds, ...)
self.setCentralWidget(self.planner)

self.termine_dock = TermineDock(self)
self.addDockWidget(Qt.LeftDockWidgetArea, self.termine_dock)
\end{verbatim}

Als Grundlage wird \texttt{QMainWindow} von Qt verwendet.\footnote{\url{https://doc.qt.io/qt-6/qmainwindow.html}} Dadurch kann die Anwendung einen zentralen Bereich und mehrere zusätzliche Dock-Bereiche besitzen. Der zentrale Bereich enthält den Planner Workspace mit den Kalenderansichten. Die Dock-Bereiche enthalten die Terminliste, den Dateneditor, die Filter, die Datumsnavigation und die Konfliktanzeige.

Im Hauptfenster werden auch die Menüs und andere Aktionen, wie Aktualisieren, Einstellungen, Undo, Importieren und Exportieren angelegt. Diese Aktionen werden mit den jeweiligen Funktionen der Anwendung verbunden. Dadurch können wichtige Arbeitsschritte direkt über die Menüleiste oder über Tastenkombinationen ausgeführt werden.

Die Verbindung zwischen den einzelnen Bereichen erfolgt über das Signal-System von Qt.\footnote{\url{https://doc.qt.io/qt-6/signalsandslots.html}} Dadurch können Aktionen aus einem Bereich an andere Bereiche der Anwendung weitergegeben werden. Ein Beispiel dafür ist die Verbindung zwischen Konfliktanzeige und dem Kalender durch Highlights:

\begin{verbatim}
self.conflicts_dock.conflict_items_highlight.connect(self.planner.highlight_termine)
\end{verbatim}^

Wird ein Konflikt ausgewählt, können die betroffenen Termine im Kalender hervorgehoben werden.

Eine weitere wichtige Aufgabe des Hauptfensters ist die Aktualisierung der Benutzeroberfläche. Wenn ein Termin erstellt, bearbeitet, gelöscht oder verschoben wird, müssen mehrere Bereiche den neuen Stand anzeigen. Dafür werden zentrale Aktualisierungsmethoden verwendet wie:

\begin{verbatim}
def refresh_everything(self) -> None:
    self.planner.refresh(emit=False)
    self.refresh_docks()
\end{verbatim}

Dabei wird zuerst der Planner Workspace aktualisiert. Anschließend werden über \texttt{refresh\_docks()} die weiteren Bereiche wie Terminliste, Filter, Dateneditor und Konfliktanzeige neu geladen.

Auch Funktionen, die nicht direkt zur Terminplanung gehören, werden im Hauptfenster eingebunden. Dazu zählen zum Beispiel die Layoutverwaltung und die Tastenkombinationen. Der \texttt{LayoutManager} speichert und lädt verschiedene Anordnungen der Oberfläche. Die Tastenkombinationen werden in einer eigenen Datei gesammelt und beim Start des Hauptfensters registriert:

\begin{verbatim}
self.layout_mgr = LayoutManager(self)
install_main_window_shortcuts(self)
\end{verbatim}

Häufige Aktionen wie Aktualisieren, Ansichtswechsel oder das Erstellen neuer Termine lassen sich dadurch schneller ausführen. Außerdem können eigene Layouts der Benutzeroberfläche gespeichert und verwaltet werden.

\subsection{Kalenderansichten}

\subsection{Termin- und Stammdatenverwaltung}

\subsection{Drag-and-Drop-Funktionalität}

\subsection{Filterfunktion}

\subsection{Konflikterkennung}

\subsection{Import, Export und Speicherung}

\subsection{Undo/Redo, Shortcuts und Layout-Verwaltung}

\newpage

\section{Test und Evaluation}

\subsection{Teststrategie}

\subsection{Funktionale Tests}

\subsection{Tests anhand typischer Planungsszenarien}

\subsection{Bewertung der Ergebnisse}

\newpage

\section{Diskussion}

\subsection{Zielerreichung}

\subsection{Stärken der entwickelten Lösung}

\subsection{Grenzen der Software}

\subsection{Erweiterungsmöglichkeiten}

\newpage

\section{Fazit und Ausblick}

\newpage

\section{Literaturverzeichnis}

\end{document}